\documentclass[10pt,a4paper]{amsart}
\usepackage[margin=19mm]{geometry}
\usepackage{amsmath,amssymb,mathtools}
\usepackage[scr=rsfs,cal=euler]{mathalfa}
\usepackage{xfrac}
\usepackage[unicode,hidelinks,bookmarks=false]{hyperref}
\newcommand{\ie}{i.e.\ }
\newcommand{\cf}{cf.\ }
\newcommand{\resp}{respectively\ }
\newcommand{\Subsection}{\subsection}
\providecommand{\nobreakdash}{\nobreak\hbox{-}}
\setlength{\emergencystretch}{2em}
\multlinegap0pt
\theoremstyle{plain}
\newtheorem{theorem}{Theorem}
\title[The growth rate of the expected number of identifiable components]{Finite Admixture Models: a Bridge with Stochastic Geometry and Choquet Theory: the growth rate of the expected number of identifiable admixture components (Theorem 3)}
\author{Michele Caprio, Sayan Mukherjee}
\date{Source: arXiv:2002.08409v12 (CC BY 4.0)}
\begin{document}
\begin{abstract}
Given a finite admixture model whose components and weights are unknown, let the number of identifiable components be a function of the amount of data sampled from a known distribution on the unit simplex. We use techniques from stochastic convex geometry to find the growth rate of its expected value. In addition, when the components are known but the weights are not, we provide an application of the classic Glivenko-Cantelli's theorem that allows us to retrieve the Choquet measure supported on the identifiable admixture components. In turn, this gives us the identifiable admixture weights. Finally, we propose a novel algorithm that estimates the model capturing the complexity of the data using only the strictly necessary number of components.


\end{abstract}
\maketitle

\noindent\emph{Source and licence.} Finite Admixture Models: a Bridge with Stochastic Geometry and Choquet Theory. Version arXiv:2002.08409v12 (CC BY 4.0), distributed under the Creative Commons Attribution 4.0 International licence, http://creativecommons.org/licenses/by/4.0/, as recorded in the arXiv licence metadata for this exact version and shown on the abstract page https://arxiv.org/abs/2002.08409v12. Full rights notice: licenses/arxiv-2002.08409-CC-BY-4.0.txt.\par\smallskip

\noindent\textit{Editorial scope.} Counted unit: the growth-rate theorem for the expected number of identifiable admixture components, printed as Theorem 3 of the article (source0.tex lines 421--433, source label \texttt{growth}), delivered together with the article's complete original proof of it, which the article places in its appendix (source0.tex lines 1909--1913, one \texttt{proof} environment carrying the article's own optional heading ``Proof of Theorem \texttt{\textbackslash ref\{growth\}}'', five lines, with no gap and no deferral). The counted statement is the article's first main result and is announced as such in its abstract and in its section on main results. No mathematics is written, repaired or re-derived: statement and proof are the article's own text line for line, in the article's own notation ($J$, $n$, $S_1,\dots,S_n$, $K_n$, $M(n)$, $F_0(K_n)$, $\Delta^{J-1}$, $\tau$, $c(J)$, $\mathbb{E}$, $\mathbb{V}$). The article prints the theorem's descriptive title as the first, bold, element of its statement body rather than as an optional theorem note, and the wrapper reproduces that construction exactly, so the delivered theorem head is followed by the article's own bold title. Required context retained uncounted: the article's abstract, reproduced verbatim as front matter (line 141); the article's own definition of the model, with the sampling display $S_1,\dots,S_n\stackrel{iid}{\sim}\mbox{Uniform}(\Delta^{J-1})$ labelled \texttt{eq1\_fmm}, the random polytope $K_n$, and the definition of $M(\cdot)$ as the cardinality of the extremal set of $K_n$ (lines 377--382); the article's own geometric preliminaries defining $i$-faces and the number $F_i(P)$ of $i$-faces of a polytope (lines 411--414) together with the sentence that identifies $F_0(K_n)$ with the number of extremal elements and so with $M(n)$ (line 416); the article's sentence introducing the counted theorem (line 419); and the article's own definition of a tower, or flag, of a polytope, denoted $\tau(P)$ (line 1907), which is the definition the retained proof uses in its last step. Every definition, numbered display, label and structural statement of those lines is retained, unchanged. Omitted from this package and not redistributed: the article's preamble and titlepage block (its authors, addresses and e-mail addresses, its keywords and subject classification, and its own macro and theorem declarations), its whole introduction and literature review, its sections on the setup of the model, on Choquet theory, on the P\'olya tree approximation, on population genetics and its conclusions, the remainder of its appendix, and its bibliography other than the four entries transcribed below; in the source's own line numbering those are lines 1--140, and the lines between and after the retained excerpts, namely 142--376, 383--410, 415, 417--418, 420, 434--1906, 1908, and 1914--2131. In particular none of the article's seven figure environments is retained. Nothing inside a retained excerpt is omitted or altered; the unit's own text outside the excerpts is the wrapper matter described here. The article's own proof of its Theorem 4, of its Theorem 5 and of its other statements are not retained, and the article's proof of its Theorem 3 is delivered in full. Citations: the retained excerpts carry six active citation commands, which cite exactly four keys of the article's own bibliography, \texttt{vu}, \texttt{ziegler}, \texttt{reitzner} and \texttt{barany2}. Those four entries are transcribed from the article's own external bibliography file that ships with the e-print (the 47-entry \texttt{main.bbl}, whose entries the article prints under the \texttt{plain} style) with the article's own printed numbers, so the delivered citation marks read as the article's do (\texttt{[6]}, \texttt{[40]}, \texttt{[43]} and \texttt{[47]} with the article's own optional notes). The only change in the transcription is that BibTeX's \texttt{\textbackslash newblock} line-break instruction, which is a typesetting directive and not part of the reference text, is removed and the entry's text is joined into a single paragraph with the article's own wording, italics and name accents kept. No other entry of the article's bibliography is transcribed or redistributed. Reference adaptation: none was needed and none was made. Every reference inside the delivered unit resolves inside it: the label \texttt{eq1\_fmm} of the retained sampling display, which the counted statement cites; the labels \texttt{first\_theorem} and \texttt{equation\_imp\_first} of the two numbered displays inside the counted statement, which the article's retained introducing sentence and the counted proof both cite; and the label \texttt{growth} of the counted statement itself, which the article's own optional proof heading cites. The label \texttt{growth} is the article's own anchor. No unresolved reference or citation is produced. Printed slips kept literal: no printed word, symbol, hypothesis or number of the counted statement or of its proof was altered. In particular the article's own en dash in the phrase ``Bachmann--Landau big-O notation'' inside the counted statement is the article's own character, delivered as it stands; the article's own spelling of its sampling display with the superscript $iid$ over a relation symbol is kept; the proof's own reference to the constant $\tau(\Delta^{J-1})=(J-1)!\cdot J$ is kept as printed; and the article's own phrase ``follow immediately after realizing that'' is kept. Layout and class: the article's own class is \texttt{amsart} with the options \texttt{12pt} and \texttt{reqno}, the same class family as the wrapper but at 12pt where the wrapper is pinned at 10pt on A4, and no class or style file travels with this package because the article uses none beyond \texttt{amsart} itself. The article's own preamble loads \texttt{fontenc}, \texttt{algorithm}, \texttt{algpseudocode}, \texttt{bbm}, \texttt{tikz}, \texttt{fancyhdr} and others; no retained excerpt needs any of them, so the wrapper loads none of that machinery. The wrapper declares the theorem environment the counted statement needs and adds no author macro, because every control sequence used by the retained excerpts is standard (the script alphabet written with the article's \texttt{\textbackslash mathscr} command, which the wrapper's pinned \texttt{mathalfa} already provides, is the only non-default one). The article numbers its theorem-like environments in one global counter, and the wrapper does the same, so the counted statement prints as Theorem 1 of the unit where the article prints it as Theorem 3. The article's equation counter is global and its displayed equations run through the whole paper, so the retained sampling display prints as (1) of the unit where the article prints it as its equation (3), and the two displays inside the counted statement print as (2) and (3) of the unit where the article prints them as its equations (4) and (5); the wrapper's retained copy of the counted proof therefore cites (2) and (3) where the article's own text cites (4) and (5). The 12pt-to-10pt difference is a page-appearance difference of the wrapper only, and no formula, symbol, hypothesis, constant or argument changes. Assets: the article has thirteen graphics-inclusion commands and seven figure environments, and six of the thirteen targets (\texttt{algo1}, \texttt{algo2}, \texttt{algo3}, \texttt{algo4}, \texttt{frac\_1} and \texttt{frac\_2}) have no file in the e-print at all; none of those items lies in a retained excerpt, no retained excerpt contains a graphics-inclusion command, a PDF-page inclusion or an \texttt{input} directive, the package declares no asset dependency, and no image file is copied into or redistributed with it. All mathematics of the unit is native text with no image dependency, so the absent graphics targets of the article raise no question for this delivery. Licence: version arXiv:2002.08409v12 of this article is distributed under the Creative Commons Attribution 4.0 International licence, http://creativecommons.org/licenses/by/4.0/, as recorded in the arXiv licence metadata for this exact version and shown on the abstract page https://arxiv.org/abs/2002.08409v12. A full rights notice is delivered at licenses/arxiv-2002.08409-CC-BY-4.0.txt. Characters: the retained excerpts contain exactly one character outside ASCII, the article's own en dash in the phrase Bachmann--Landau inside the counted statement; the delivered engine is XeLaTeX, whose default Latin Modern text font carries it, and the compile receipt reports no missing character. No glyph is substituted and no word is rewritten.
Suppose the number $M$ of identifiable admixture components is a function $M(n)$ of the amount $n$ of data $x_1,\ldots,x_n$ that we sample. Such function is defined as follows. Let
\begin{equation}\label{eq1_fmm}
S_1,\ldots,S_n \stackrel{iid}{\sim} \mbox{Uniform}(\Delta^{J-1}), \quad J\geq 2,
\end{equation}
and call $K_n:=\text{Conv}(s_1,\ldots,s_n)$, where $s_j$ denotes the realization of $S_j$. In the stochastic convex geometry literature \cite{vu}, $K_n$ is called a \textit{random polytope}. Function $M(\cdot)$ is defined as
$$M:\mathbb{N} \rightarrow \mathbb{N}, \quad n \mapsto M(n):=\#\text{ex}(K_n),$$

\begin{itemize}
    \item As pointed out in \cite[Definition 2.1]{ziegler}, in higher-dimensional geometry, the faces of a polytope are features of all dimensions. A face of dimension $i$ is called an \textit{$i$-face}. For example, the polygonal faces of an ordinary polyhedron are $2$-faces. For any $h$-dimensional polytope, we have that $-1 \leq i \leq h$, where $-1$ is the dimension of the empty set. Let us  give a clarifying example. The faces of a cube comprise the cube itself ($3$-face), its facets ($2$-faces), the edges ($1$-faces), its vertices ($0$-faces), and the empty set (having dimension $-1$). Given a generic $h$-dimensional polytope $P$, we denote by $\mathcal{F}_i(P)$ one of its $i$-faces, $i \in \{-1,0,\ldots,h\}$.
    \item We call $\mathscr{F}_i(P)$ the \textit{collection of its $i$-faces}, and ${F}_i(P)$ the \textit{number of its $i$-faces}, that is, ${F}_i(P)=\#\mathscr{F}_i(P)$, for all $i$.
\end{itemize}

In view of the above definition of $0$-faces, we denote by $F_0(K_n)$ the number of extremal elements of $K_n$, so $M(n)=F_0(K_n)$. In the remainder of this section, we keep both notations to highlight the relationship between (stochastic convex) geometry and finite admixture models.

First we find the expected number of identifiable components and we show that it grows at rate $(\log n)^{J-1}$. Equations \eqref{first_theorem} and \eqref{equation_imp_first} are a consequence of \cite[Theorem 6]{reitzner} and \cite[Theorem 5]{barany2}, respectively.

\begin{theorem}{\textbf{(Growth rate of $\mathbb{E}[M(n)]$).}}\label{growth}
Let $K_n :=\text{Conv}({s}_1,\ldots,{s}_{n})$, where ${S}_1,\ldots,{S}_{n}$ are sampled as in \eqref{eq1_fmm}. Then,
\begin{align}\label{first_theorem}
\mathbb{E}[M(n)] = \mathbb{E} \left[F_0(K_n)\right] = \frac{J}{(J+1)^{J-1}}  (\log n)^{J-1} + \mathcal{O}\left((\log n)^{J-2}\log\log n\right),
\end{align}
where $\mathcal{O}$ denotes Bachmann–Landau big-O notation.
%and $\tau(\Delta^{J-1})$ is the number of towers of $\Delta^{J-1}$. 
In turn, this implies that
\begin{equation}\label{equation_imp_first}
    \lim_{n \rightarrow \infty} (\log n)^{-(J-1)}\mathbb{E} \left[M(n)\right] = \lim_{n \rightarrow \infty} (\log n)^{-(J-1)}\mathbb{E} \left[F_0(K_n)\right] = \frac{J}{(J+1)^{J-1}}  =: c(J).
\end{equation}
%where $T(\Delta^{J-1})$ is the number of towers of $\Delta^{J-1}$.
\end{theorem}

Recall that a chain $\mathcal{F}_0(P)\subset \mathcal{F}_1(P) \subset \cdots \subset \mathcal{F}_{h}(P)$ of $i$-dimensional faces of a polytope $P$ is called a \textit{tower} (or a \textit{flag}) of $P$, and we denote it as $\tau(P)$.

\begin{proof}[Proof of Theorem \ref{growth}]
 In \cite[Theorem 6]{reitzner} and \cite[Theorem 5]{barany2}, the authors show that, given a convex polytope $P$ in $\mathbb{R}^d$, if we call $P_n$ the convex hull of $n$ points sampled iid from a uniform on $P$, then 
$$\mathbb{E} \left[F_0(P_n)\right] = \frac{\tau(P)}{(d+1)^{d-1}  (d-1)!}  (\log n)^{d-1} + \mathcal{O} \left( (\log n)^{(d-2)} \log\log n \right).$$
Then, since $\Delta^{J-1}$ is a convex polytope in $\mathbb{R}^J$, and given the way we defined $K_n$, equations \eqref{first_theorem} and \eqref{equation_imp_first} follow immediately after realizing that $\tau(\Delta^{J-1})=(J-1)! \cdot J$.
\end{proof}

\begin{thebibliography}{99}
\bibitem[6]{barany2} Imre B{\'a}r{\'a}ny and Christian Buchta. Random polytopes in a convex polytope, independence of shape, and concentration of vertices. {\em Mathematische Annalen}, 297:467--497, 1993.
\bibitem[40]{reitzner} Matthias Reitzner. The combinatorial structure of random polytopes. {\em Advances in Mathematics}, 191(1):178--208, 2005.
\bibitem[43]{vu} Van~Hanh Vu. Sharp concentration of random polytopes. {\em Geometric and Functional Analysis}, 15:1284--1318, 2005.
\bibitem[47]{ziegler} G{\"u}nter~M. Ziegler. {\em Lectures on Polytopes}, volume 152 of {\em Graduate Texts in Mathematics}. Springer, New York, 1995.
\end{thebibliography}
\end{document}
